% !PRINTED 28
\weberpage{28}

The product of the constant factor on the left-hand side is the
discriminant $D$, and so by (13) one obtains
\weq{21}{\begin{aligned}
D\,y_1^2y_2^2y_3^2y_4^2 &= f(\Phi, -\Psi)\\
&= a_0\Phi^4 - a_1\Phi^3\Psi + a_2\Phi^2\Psi^2 - a_3\Phi\Psi^3 + a_4\Psi^4,
\end{aligned}}
and the $y_i^2$ are therefore the linear factors of this biquadratic
form.

\begin{center}
\rule{2cm}{0.4pt}
\end{center}

There is yet a further simultaneous covariant of the form $\varphi$,
$\psi$, namely the \gesperrt{Jacobian} functional determinant of the four
functions $\varphi$, $\psi$, $\Phi$, $\Psi$. We denote it thus:
\weq{22}{K = \begin{vmatrix}
\varphi_1, & \varphi_2, & \varphi_3, & \varphi_4\\
\psi_1, & \psi_2, & \psi_3, & \psi_4\\
\Phi_1, & \Phi_2, & \Phi_3, & \Phi_4\\
\Psi_1, & \Psi_2, & \Psi_3, & \Psi_4
\end{vmatrix}.}
If one forms it for the canonical form, there results, by (12),
(14), (15),
\[
K = -\Delta\,y_1y_2y_3y_4,
\]
and since $\Delta^2 = D$:
\[
K^2 = D\,y_1^2y_2^2y_3^2y_4^2,
\]
and so by (21) we have
\weq{23}{K_2 = f(\Phi, -\Psi).} % [sic] printed K_2; derivation requires K^2

\begin{center}
\rule{2cm}{0.4pt}
\end{center}

The biquadratic form $f(\Phi, -\Psi)$ has two covariants, which we
denote, as in Bd. I, §\,70, thus:
\weq{24}{\begin{aligned}
H &= \frac{1}{3}\begin{vmatrix}
\dfrac{\partial^2 f}{\partial \Phi^2}, & -\dfrac{\partial^2 f}{\partial \Phi\,\partial \Psi}\\
\dfrac{\partial^2 f}{\partial \Psi\,\partial \Phi}, & -\dfrac{\partial^2 f}{\partial \Psi^2}
\end{vmatrix},\\
T &= \frac{1}{12}\begin{vmatrix}
f'(\Phi), & f'(-\Psi)\\
H'(\Phi), & H'(-\Psi)
\end{vmatrix}.
\end{aligned}}
They are of the fourth and sixth degree in the $\Phi$, $\Psi$, and
hence of the eighth and twelfth degree in the $x_i$. They are
\gesperrt{covariants of the fundamental curve} (combinants of $\varphi$ and
$\psi$), and between them there holds the relation:
\weq{25}{H^3 - 48\,A\,H\,f^2 - 64\,B\,f^3 = -27\,T^2.}

\begin{center}
\rule{2cm}{0.4pt}
\end{center}
